\chapter{An exact Taylor language for finite choices}\label{ch:discrete-taylor}

Every finite coalition table has one exact multilinear representation. The coefficients are precisely its Möbius interactions. This gives EBU a compact mathematical language for the full hierarchy and makes clear when a reduced description is exact, rather than an approximation that discards part of the account.


\section{Why use the word Taylor?}
A continuous Taylor expansion builds a function from its value, slopes, curvature and higher derivatives near a point. A coalition table has a related hierarchy: its baseline, single-action changes, pair interactions, and higher interactions. The resemblance becomes an exact mathematical statement when we describe absent/present choices by indicator variables.

There is also a crucial difference. An ordinary truncated continuous Taylor series generally leaves a remainder. The complete multilinear representation of a finite Boolean table leaves none. It reproduces every row exactly because it has exactly the coefficients needed to represent that finite collection of choices.

\section{From a subset to its indicators}
Let $z_i$ be one when action $i$ is included and zero otherwise. Write $\chi_S$ for the vector representing $S$. Define
\begin{equation}\label{eq:multilinear}
 \widetilde E(z)=\sum_{T\subseteq N}m_E(T)\prod_{i\in T}z_i,
 \qquad z\in\mathbb R^n,
\end{equation}
with the empty product equal to one. Each variable appears with exponent at most one, which is what multilinear means here.

\begin{theorem}{Exact finite discrete Taylor representation}\label{thm:multilinear}
Equation~\ref{eq:multilinear} is the unique multilinear polynomial satisfying $\widetilde E(\chi_S)=E(S)$ for all $S\subseteq N$. It has no remainder on the Boolean cube, and
\begin{equation}
 \partial_T\widetilde E(0)=m_E(T)
\end{equation}
for each set of distinct variables $T$.
\end{theorem}
\begin{proof}
At $\chi_S$, a product is one precisely when $T\subseteq S$ and otherwise zero. Evaluation therefore gives $\sum_{T\subseteq S}m_E(T)=E(S)$. Any multilinear polynomial has one coefficient per such product; evaluating its subsets in increasing size determines these coefficients uniquely by the same inversion. Differentiating once in each variable of $T$ and evaluating all variables at zero selects exactly its coefficient.
\end{proof}

The theorem allows the whole landscape to be reconstructed from its interactions. It does not require small physical actions, differentiable endpoints or an equilibrium probability law. It is a representation theorem for the finite table.

\section{A complete quadratic example}
For the three equal one-unit couriers of the previous chapter,
\begin{equation}
 \widetilde E(z)=\frac{15}{4}(z_1+z_2+z_3)
             -\frac12(z_1z_2+z_1z_3+z_2z_3).
\end{equation}
There is no $z_1z_2z_3$ term. Setting all indicators to one gives $45/4-3/2=39/4$, exactly the triple value. Setting only $z_1,z_2$ to one gives $15/2-1/2=7$, exactly the pair value.

Now compare this with the physical polynomial for fractional control intensities,
\begin{equation}
 E_{\mathrm{phys}}(z)=4(z_1+z_2+z_3)-\tfrac14(z_1+z_2+z_3)^2.
\end{equation}
At Boolean vertices $z_i^2=z_i$, so it reduces to the multilinear formula above. Between vertices the two functions generally differ. For one half-active action, the physical polynomial gives $31/16$, whereas the multilinear extension gives $15/8$. Both agree when the action is entirely absent or present.

This small difference carries a large conceptual lesson: the unique multilinear extension is not automatically the physical law for fractional actions. Binary observations determine one interpolating polynomial in the multilinear class. They do not determine every smooth function through the same vertices.

\section{A probabilistic interpolation with a declared meaning}
There is a useful alternative interpretation. Suppose actions are included independently, with inclusion probabilities $z_i\in[0,1]$, and then the selected coalition receives its table value. The expected table value is
\begin{equation}
 \sum_{S\subseteq N}E(S)
 \prod_{i\in S}z_i\prod_{j\notin S}(1-z_j)
 =\widetilde E(z).
\end{equation}
The equality follows by taking expectations of the product indicators in Equation~\ref{eq:multilinear}. This is a mathematical randomized selection model with independent inclusions. It does not claim that actual actors choose independently or that an infeasible group becomes admissible because it has a probability weight.

A fractional physical dose and a probability of a full dose are different interventions. Confusing them can turn an exact interpolation into a false physical prediction. In an EBU system, the record should say whether a slider changes quantity, participation probability, or merely the displayed interpolation.

\section{Constants, units and reparametrization}
The empty coefficient records the baseline table value. In the normalized action table it is zero. With a natural-drift comparator it may not be, and we must retain it or explicitly center the table. Nonempty differences annihilate a constant, but reconstruction of the full value still needs that constant.

The degree also refers to declared action coordinates. Splitting an action into new controls, changing a resolver or choosing nonlinear intensity coordinates can change the apparent polynomial structure of the physical composite. A representation theorem does not make every coordinate choice equally meaningful. Physical interpretation requires the action labels and response map that supplied the table.

\section{Why this helps an EBU system}
The complete polynomial is a second representation of the same table. It can reconstruct every coalition value from the baseline and interaction coefficients. This supports explanation and reconciliation: a reader can see which single actions, pairs and larger groups contribute to a proposed total.

If a theorem limits the degree, the representation becomes genuinely smaller. A quadratic potential with affine endpoint response needs only the baseline, singleton and pair coefficients. A common linear feedback system preserves that ceiling even if it contains many temporal modes. The simplification follows from the composed function's degree, not from the number of differential equations or the visual smoothness of a trajectory.

Without such a structural result, deleting higher orders changes the table. The omitted terms have an exact remainder on each coalition: their sum over subsets of that coalition. This lets an EBU design state the consequence of a reduced account explicitly instead of hiding it in a convenient approximation.

The efficiency benefit is precise. Use the smallest representation that retains the required information, and preserve the full dependencies where they matter. This reduces unnecessary computation while helping a community understand why a coordinated service has its reported value.

\section{The boundary with ordinary calculus}
The word exact belongs to the finite Boolean representation. The next chapter connects a fixed additive physical response to mixed derivatives of the state potential. That connection explains why small interactions have familiar Taylor orders and why quadratic fields give pair-only structure under affine action maps.

The two constructions will meet, but they will not become interchangeable. The discrete polynomial is determined by all vertices. The local continuous derivative is determined near a point. One is a finite table identity; the other is local geometric information with explicit regularity and remainder conditions.
